\chapter{Poincaré 对称性与 Noether 定理}

物理学的核心信条之一是：对称性决定守恒律。一个系统如果在空间平移下不变，就拥有动量守恒；如果在时间平移下不变，就拥有能量守恒；如果在旋转下不变，就拥有角动量守恒。Noether 定理把这一直觉提升为精确的数学命题：每一个连续对称性对应一个守恒流。在非相对论力学中，这只涉及 Galilei 群的三维平移和三维旋转。一旦进入相对论，时空统一为四维 Minkowski 流形，对称群变为 Poincaré 群 $ISO(1,d-1)$，它同时包含 Lorentz 变换（旋转和 推进）与四维平移。Poincaré 群的结构比 Galilei 群更丰富，也比纯旋转群更复杂——它是一个半直积，平移子群与 Lorentz 子群之间有非平凡的交织。

前十五章建立了群、表示论、Clifford 代数与微分几何的语言。本章把这些工具汇合于一个物理问题：Poincaré 群如何支配平直时空中的场论。路线图是这样的。首先从 Poincaré 群的群乘法读出它的半直积结构，再通过 Lie 代数把有限变换的无穷小生成元 $P^\mu$（平移）和 $M^{\mu\nu}$（Lorentz）的交换关系一次性导出。然后问：一个物理场 $\Phi^A(x)$ 在 Poincaré 变换下如何变化？答案把场的变换分成两部分——时空坐标移动带来的"轨道"部分和场分量本身在内部表示中旋转的"内禀"部分。有了场的变换律，就可以把它代入作用量变分，Noether 定理自动给出能动量张量和角动量流的守恒方程。最后，Belinfante--Rosenfeld 改进把自旋流吸收进对称的能动量张量，为下一章的曲率时空推广做好准备。

\section{记号与约定}

在进入群结构之前，先固定整套记号。度规符号的选取、生成元的归一化以及无穷小展开的约定，会在后续每一步推导中反复出现；只要其中一项前后不一致，结构常数的符号或 $i$ 因子就会出错。

在 $d$ 维 Minkowski 时空中取度量
\begin{equation*}\eta_{\mu\nu}=\operatorname{diag}(-,+,\ldots,+),\qquad \mu,\nu=0,\ldots,d-1.
\end{equation*}
这里选"多为正"符号约定，即时间分量负、空间分量正。这个选取在广义相对论和粒子物理中都很常见；它的好处是空间度规正定，三维 Euclidean 几何的公式可以直接继承。升降指标均由 $\eta_{\mu\nu}$ 完成，例如 $V^\mu=\eta^{\mu\nu}V_\nu$，$V_0=-V^0$，$V_i=V^i$。全文使用自然单位 $\hbar=c=1$。

Poincaré 群在量子态空间上的表示是酉算符。无穷小展开约定为
\begin{equation}\label{eq:ms-poincare-U-expand}
\boxed{
U(1+\omega,\epsilon)=I+i\epsilon_\mu P^\mu-\frac{i}{2}\omega_{\mu\nu}M^{\mu\nu}+O(\epsilon^2,\omega^2),
}
\end{equation}
其中 $P^\mu$ 是平移生成元（物理上是四动量算符），$M^{\mu\nu}=-M^{\nu\mu}$ 是 Lorentz 生成元（物理上是角动量算符，包含空间旋转和 推进）。三个细节值得解释。第一，因子 $i$ 使 $P^\mu$ 和 $M^{\mu\nu}$ 成为 Hermite 算符，从而对应物理可观测量。第二，$\omega_{\mu\nu}$ 的反对称性 $\omega_{\mu\nu}=-\omega_{\nu\mu}$ 保证只有 $d(d-1)/2$ 个独立 Lorentz 参数（$d=4$ 时是 6 个：3 个旋转加 3 个 推进），而求和中的 $1/2$ 补偿了每对 $(\mu,\nu)$ 被计算两次。第三，平移参数 $\epsilon^\mu$ 是一个 $d$-矢量，没有反对称约束。

\section{Poincaré 群及其 Lie 代数}
物理上，Poincaré 变换是保持 Minkowski 间隔 $\dd s^2=\eta_{\mu\nu}\dd x^\mu \dd x^\nu$ 不变的最一般仿射变换。它由一个 Lorentz 矩阵 $\Lambda$ 和一个平移矢量 $a$ 组成：
\begin{equation*}x'^\mu=\Lambda^\mu{}_\nu x^\nu+a^\mu.
\end{equation*}
Lorentz 矩阵必须满足
\begin{equation}\label{eq:ms-poincare-lorentz-def}
\boxed{\Lambda^T\eta\Lambda=\eta,}
\end{equation}
即 $\eta_{\rho\sigma}\Lambda^\rho{}_\mu\Lambda^\sigma{}_\nu=\eta_{\mu\nu}$。这个条件的几何意义是：$\Lambda$ 保持 Minkowski 内积不变。物理上，它保证光锥结构不变——如果两个事件由光信号连接（$\dd s^2=0$），变换后仍然如此。
群乘法告诉我们连续做两次 Poincaré 变换的净效果。设 $g_1=(\Lambda_1,a_1)$ 和 $g_2=(\Lambda_2,a_2)$。先做 $g_1$：$x'=\Lambda_1 x+a_1$；再做 $g_2$：$x''=\Lambda_2 x'+a_2$。代入得
\begin{equation}
x''=\Lambda_2\Lambda_1 x+(\Lambda_2 a_1+a_2).
\end{equation}
因此群乘法为
\begin{equation}\label{eq:ms-poincare-semidirect}
\boxed{
(\Lambda_2,a_2)(\Lambda_1,a_1)=(\Lambda_2\Lambda_1,\;\Lambda_2 a_1+a_2).
}
\end{equation}
关键观察：平移部分出现的是 $\Lambda_2 a_1$ 而非 $a_1$。这意味着先平移再 Lorentz 变换，与先 Lorentz 变换再平移，效果不同——Lorentz 变换会"搅动"平移矢量。这正是半直积的特征：
\begin{equation}
ISO(1,d-1)=SO(1,d-1)\ltimes\mathbb{R}^{1,d-1}.
\end{equation}
如果是直积，群乘法会是 $(\Lambda_2\Lambda_1,\;a_2+a_1)$，两个子群完全独立。半直积中的 $\ltimes$ 记号正是记录了 Lorentz 对平移的非平凡作用。这个结构会在下一步直接进入 Lie 代数的交换关系：$[M,P]\neq 0$。
逆元公式在推导 Lie 括号 时会用到。令 $(\Lambda,a)^{-1}=(\Lambda',a')$，要求 $(\Lambda',a')(\Lambda,a)=(I,0)$。由群乘法\eqnrefs{eq:ms-poincare-semidirect} 得 $(\Lambda'\Lambda,\;\Lambda'a+a')=(I,0)$。第一分量给出 $\Lambda'=\Lambda^{-1}$；代入第二分量给出 $a'=-\Lambda^{-1}a$。因此
\begin{equation}\label{eq:ms-poincare-inverse}
(\Lambda,a)^{-1}=(\Lambda^{-1},\;-\Lambda^{-1}a).
\end{equation}
注意平移部分的逆元不是 $-a$ 而是 $-\Lambda^{-1}a$——又是半直积的体现。

\section{Poincaré 代数的统一推导}

有了群结构，下一步问：为什么要把群变成 Lie 代数？原因是群乘法是非线性的，而 Lie 代数的交换关系是线性的结构常数，后者远比前者容易操作。无穷小生成元的交换关系 $[P,P]$、$[M,P]$、$[M,M]$ 完全编码了群的局域结构，而且可以用来构造 Casimir 不变量、分类不可约表示、验证场方程的协变性。本节不直接列出这三组公式，而是从群乘法和 Lie 括号 的定义出发，用一个统一的推导一次性得到全部三组关系。
先单独看 Lorentz 部分。取恒等元附近 $\Lambda=I+\omega+O(\omega^2)$，其中 $\omega$ 是无穷小矩阵。代入 Lorentz 条件\eqnrefs{eq:ms-poincare-lorentz-def}：
\begin{equation*}
(I+\omega)^T\eta(I+\omega)=\eta.
\end{equation*}
展开到一阶：
\begin{equation*}
\eta+\omega^T\eta+\eta\omega+O(\omega^2)=\eta,
\end{equation*}
故
\begin{equation}
\omega^T\eta+\eta\omega=0.
\end{equation}
这个条件说 $\omega$ 关于 $\eta$ 反对称。为了把矩阵条件转成分量条件，定义降指标版本 $\omega_{\mu\nu}=\eta_{\mu\rho}\omega^\rho{}_\nu$。那么
\begin{equation}
\omega_{\nu\mu}=\eta_{\nu\rho}\omega^\rho{}_\mu=\eta_{\nu\rho}\eta_{\mu\sigma}\omega^\sigma{}_\nu\cdot(\text{转置})=-\eta_{\mu\rho}\omega^\rho{}_\nu=-\omega_{\mu\nu},
\end{equation}
即
\begin{equation*}\boxed{\omega_{\mu\nu}=-\omega_{\nu\mu}.}
\end{equation*}
$d\times d$ 反对称矩阵有 $d(d-1)/2$ 个独立分量，这就是 Lorentz 群的维数。在 $d=4$ 时，6 个生成元分为 3 个旋转 $J_i=\frac{1}{2}\epsilon_{ijk}M^{jk}$ 和 3 个 推进 $K_i=M^{0i}$。
现在进入核心推导。Poincaré Lie 代数的元素写成 $X=(\omega,a)$，其中 $\omega\in\mathfrak{so}(1,d-1)$ 是反对称矩阵，$a\in\mathbb{R}^{1,d-1}$ 是平移矢量。Lie 括号 的一般定义是通过群交换子：
\begin{equation}\label{eq:ms-poincare-bracket-def}
g_1(t)g_2(s)g_1(t)^{-1}g_2(s)^{-1}=I+ts\,[X_1,X_2]+O(t^2,s^2,t^2s,ts^2),
\end{equation}
其中 $g_1(t)=(e^{t\omega_1},ta_1)$ 和 $g_2(s)=(e^{s\omega_2},sa_2)$ 是两条一参数子群。群交换子的几何意义是：先沿 $g_1$ 走一小段 $t$，再沿 $g_2$ 走一小段 $s$，然后沿 $g_1$ 退回，再沿 $g_2$ 退回。如果群是 Abel 的，这四步完全抵消；非抵消的部分到 $ts$ 阶正是 Lie 括号，它测量两个生成元"不交换"的程度。

\begin{derivation}[从群乘法导出半直积 Lie 括号]
逐次使用群乘法\eqnrefs{eq:ms-poincare-semidirect} 与逆元公\eqnrefs{eq:ms-poincare-inverse}，把四元群乘积\eqnrefs{eq:ms-poincare-bracket-def} 展开到 $ts$ 阶。关键是跟踪每一步中 Lorentz 部分如何作用于平移部分。

计算 $g_1(t)g_2(s)$。用群乘法，Lorentz 部分相乘、平移部分被第一个 Lorentz 作用后再加：
\begin{equation*}
g_1(t)g_2(s)=\big(e^{t\omega_1}e^{s\omega_2},\;e^{t\omega_1}(sa_2)+ta_1\big)+O(ts^2,t^2s).
\end{equation*}
这里 $e^{t\omega_1}(sa_2)$ 就是半直积的体现：$g_1$ 的 Lorentz 部分作用到了 $g_2$ 的平移部分上。

由逆元公\eqnrefs{eq:ms-poincare-inverse}，
\begin{equation*}
g_1(t)^{-1}=\big(e^{-t\omega_1},\;-e^{-t\omega_1}(ta_1)\big)+O(t^2).
\end{equation*}

把上述结果代入群乘法，计算 $g_1(t)g_2(s)g_1(t)^{-1}$。其 Lorentz 部分为
\begin{equation*}
e^{t\omega_1}e^{s\omega_2}\cdot e^{-t\omega_1}.
\end{equation*}
平移部分（第一分量的 Lorentz 作用于第二分量的平移，再加自身平移）：
\begin{equation*}
e^{t\omega_1}e^{s\omega_2}\big(-e^{-t\omega_1}ta_1\big)+\big(e^{t\omega_1}sa_2+ta_1\big).
\end{equation*}
合并写为
\begin{equation*}
g_1g_2g_1^{-1}=\Big(e^{t\omega_1}e^{s\omega_2}e^{-t\omega_1},\;\;-e^{t\omega_1}e^{s\omega_2}e^{-t\omega_1}ta_1+e^{t\omega_1}sa_2+ta_1\Big).
\end{equation*}

展开到 $ts$ 阶。Lorentz 部分：利用 Baker--Campbell--Hausdorff 或直接展开矩阵指数，
\begin{equation*}
e^{t\omega_1}e^{s\omega_2}e^{-t\omega_1}=(I+t\omega_1)(I+s\omega_2)(I-t\omega_1)+O(t^2,s^2)=I+s\omega_2+O(s^2).
\end{equation*}
$t$ 的一阶项被共轭消去了——这是预期的，因为共轭 $\Lambda\omega_2\Lambda^{-1}$ 到 $t=0$ 就是 $\omega_2$ 本身。$ts$ 阶的修正要到随后右乘 $e^{-s\omega_2}$ 时才出现。

平移部分：把指数展开到一阶，
\begin{align*}
&-e^{t\omega_1}e^{s\omega_2}e^{-t\omega_1}ta_1+e^{t\omega_1}sa_2+ta_1\\
&\approx -(I+t\omega_1+s\omega_2)(I-t\omega_1)ta_1+(I+t\omega_1)sa_2+ta_1\\
&= -(I-t\omega_1+t\omega_1-t^2\omega_1^2+s\omega_2-ts\omega_2\omega_1)ta_1+ts\omega_1 a_2+ta_1+O(t^2,ts^2)\\
&= -ta_1+t^2\omega_1^2 a_1-ts\omega_2 a_1+ts\omega_1 a_2+ta_1+O(t^2,ts^2)\\
&= ts(\omega_1 a_2-\omega_2 a_1)+O(t^2,ts^2).
\end{align*}
$t^2$ 阶项略去。关键：$ta_1$ 和 $-ta_1$ 抵消，留下的是 $ts$ 阶的 $\omega_1 a_2-\omega_2 a_1$。这一项正是两个生成元"交叉作用"的产物：$\omega_1$ 作用到 $a_2$ 上，减去 $\omega_2$ 作用到 $a_1$ 上。

最后右乘 $g_2(s)^{-1}=(e^{-s\omega_2},-e^{-s\omega_2}sa_2)$。将前式记为 $(I+s\omega_2,\;ts(\omega_1 a_2-\omega_2 a_1))$，其 Lorentz 部分为
\begin{equation*}
(I+s\omega_2)(I-s\omega_2)+O(s^2)=I+O(s^2).
\end{equation*}
平移部分：$(I+s\omega_2)$ 作用于 $(-e^{-s\omega_2}sa_2)$ 再加 $ts(\omega_1 a_2-\omega_2 a_1)$：
\begin{equation*}
(I+s\omega_2)(-sa_2+O(s^2))+ts(\omega_1 a_2-\omega_2 a_1)=-sa_2+ts(\omega_1 a_2-\omega_2 a_1)+O(s^2).
\end{equation*}
但还要加上 $g_2(s)^{-1}$ 自身的平移 $-e^{-s\omega_2}sa_2\approx -sa_2$，两项 $-sa_2$ 抵消，留下
\begin{equation*}
ts(\omega_1 a_2-\omega_2 a_1)+O(t^2,ts^2).
\end{equation*}

读出 Lie 括号。群交换子到 $ts$ 阶为
\begin{equation*}
g_1g_2g_1^{-1}g_2^{-1}=\Big(I+ts[\omega_1,\omega_2],\;\;ts(\omega_1 a_2-\omega_2 a_1)\Big)+O(t^2,s^2).
\end{equation*}
Lorentz 部分的 $ts$ 阶来自 $e^{t\omega_1}e^{s\omega_2}e^{-t\omega_1}e^{-s\omega_2}=I+ts[\omega_1,\omega_2]+\cdots$，这是纯 Lorentz 部分的群交换子。与\eqnrefs{eq:ms-poincare-bracket-def} 比较，读出
\begin{equation}\label{eq:ms-poincare-unified-bracket}
\boxed{
[(\omega_1,a_1),(\omega_2,a_2)]=\big([\omega_1,\omega_2],\;\omega_1 a_2-\omega_2 a_1\big).
}
\end{equation}
\end{derivation}

这条统一的 括号 公\eqnrefs{eq:ms-poincare-unified-bracket} 是整个 Poincaré 代数的压缩形式。它的两个分量分别给出：Lorentz 部分的矩阵交换子 $[\omega_1,\omega_2]$（纯 Lorentz 子代数），以及平移部分的 $\omega_1 a_2-\omega_2 a_1$（Lorentz 对平移的交叉作用）。下面三小节只需在\eqnrefs{eq:ms-poincare-unified-bracket} 中代入不同的特例，即可一次性得到全部三组交换关系。
纯平移对应 $\omega_1=\omega_2=0$。代入\eqnrefs{eq:ms-poincare-unified-bracket}：
\begin{equation*}
[(0,a_1),(0,a_2)]=\big([0,0],\;0\cdot a_2-0\cdot a_1\big)=(0,0).
\end{equation*}
两个纯平移的 括号 为零，因为 Lorentz 部分恒为零、交叉项中 $\omega=0$ 也为零。翻译成生成元语言：
\begin{equation}\label{eq:ms-poincare-PP}
\boxed{[P_\mu,P_\nu]=0.}
\end{equation}
这反映平移子群 $\mathbb{R}^{1,d-1}$ 是 Abel 群——先沿 $x^\mu$ 方向平移再沿 $x^\nu$ 方向平移，与顺序无关。
取一个纯 Lorentz 元素 $X_1=(\omega,0)$ 和一个纯平移元素 $X_2=(0,a)$。代入\eqnrefs{eq:ms-poincare-unified-bracket}：
\begin{equation*}
[(\omega,0),(0,a)]=\big([\omega,0],\;\omega a-0\cdot 0\big)=(0,\omega a).
\end{equation*}
结果是一个纯平移元素，平移矢量被 $\omega$ 作用。物理上，这说 Lorentz 变换会把一个平移矢量旋转或 推进 到另一个方向。在生成元语言中，Lorentz 生成元对平移生成元的伴随作用正是矢量表示。为了写出显式公式，需要把矩阵 $\omega$ 用反对称参数 $\omega_{\rho\sigma}$ 和矢量表示的 Lorentz 生成元 $J_V^{\rho\sigma}$ 展开。矢量表示的定义是：对任意矢量 $a$，$(\omega a)^\mu=\omega^\mu{}_\nu a^\nu$。由 $\omega^\mu{}_\nu=\eta^{\mu\alpha}\omega_{\alpha\nu}$ 和 $\omega_{\alpha\nu}$ 的反对称性，
\begin{equation*}
(\omega a)^\mu=\eta^{\mu\alpha}\omega_{\alpha\nu}a^\nu=\frac{1}{2}\omega_{\rho\sigma}\big(\eta^{\rho\mu}\delta^\sigma{}_\nu-\eta^{\sigma\mu}\delta^\rho{}_\nu\big)a^\nu.
\end{equation*}
因此矢量表示的 Lorentz 生成元为
\begin{equation}\label{eq:ms-poincare-JV}
\boxed{
(J_V^{\rho\sigma})^\mu{}_\nu=i\big(\eta^{\rho\mu}\delta^\sigma{}_\nu-\eta^{\sigma\mu}\delta^\rho{}_\nu\big),
}
\end{equation}
其中因子 $i$ 是为了与量子力学中 Hermite 生成元的约定一致。于是 $[M^{\rho\sigma},P^\mu]$ 必须按同一矢量表示作用：
\begin{equation}\label{eq:ms-poincare-MP}
\boxed{
[M^{\rho\sigma},P^\mu]=i\big(\eta^{\rho\mu}P^\sigma-\eta^{\sigma\mu}P^\rho\big).
}
\end{equation}
物理含义：Lorentz 生成元 $M^{\rho\sigma}$ 作用在平移生成元 $P^\mu$ 上，把它"旋转"到 $P^\sigma$ 和 $P^\rho$ 方向。这就是半直积在 Lie 代数中的表现——$[M,P]\neq 0$。
最后取两个纯 Lorentz 元素 $X_1=(\omega_1,0)$ 和 $X_2=(\omega_2,0)$。代入\eqnrefs{eq:ms-poincare-unified-bracket}：
\begin{equation*}
[(\omega_1,0),(\omega_2,0)]=\big([\omega_1,\omega_2],\;\omega_1\cdot 0-\omega_2\cdot 0\big)=\big([\omega_1,\omega_2],\,0\big).
\end{equation*}
平移部分为零，结果仍是纯 Lorentz 元素。Lorentz 子代数在 Poincaré 代数中是子代数（不是理想，因为 $[M,P]\neq 0$）。现在需要把矩阵交换子 $[\omega_1,\omega_2]$ 翻译成生成元 $M^{\mu\nu}$ 的交换关系。在矢量表示中，$(J^{\mu\nu})^\rho{}_\sigma=i(\eta^{\mu\rho}\delta^\nu{}_\sigma-\eta^{\nu\rho}\delta^\mu{}_\sigma)$。

\begin{derivation}[矢量表示中 Lorentz 生成元交换关系的直接计算]
先算矩阵乘积 $J^{\mu\nu}J^{\rho\sigma}$，再交换指标对 $(\mu,\nu)\leftrightarrow(\rho,\sigma)$ 得到 $J^{\rho\sigma}J^{\mu\nu}$，相减读出结构常数。计算虽然直接，但需要仔细跟踪四个 $\eta$-$\delta$ 因子的缩并。

写出矩阵乘积。对矩阵乘积 $(J^{\mu\nu}J^{\rho\sigma})^\alpha{}_\beta$，先对中间指标 $\gamma$ 求和：
\begin{equation*}
(J^{\mu\nu}J^{\rho\sigma})^\alpha{}_\beta=(J^{\mu\nu})^\alpha{}_\gamma(J^{\rho\sigma})^\gamma{}_\beta.
\end{equation*}
代入定义 $(J^{\mu\nu})^\alpha{}_\gamma=i(\eta^{\mu\alpha}\delta^\nu{}_\gamma-\eta^{\nu\alpha}\delta^\mu{}_\gamma)$ 和 $(J^{\rho\sigma})^\gamma{}_\beta=i(\eta^{\rho\gamma}\delta^\sigma{}_\beta-\eta^{\sigma\gamma}\delta^\rho{}_\beta)$：
\begin{equation*}
(J^{\mu\nu}J^{\rho\sigma})^\alpha{}_\beta=i^2\big(\eta^{\mu\alpha}\delta^\nu{}_\gamma-\eta^{\nu\alpha}\delta^\mu{}_\gamma\big)\big(\eta^{\rho\gamma}\delta^\sigma{}_\beta-\eta^{\sigma\gamma}\delta^\rho{}_\beta\big).
\end{equation*}

展开四项并对 $\gamma$ 求和。利用 $\delta^\nu{}_\gamma\eta^{\rho\gamma}=\eta^{\rho\nu}$ 等：
\begin{equation*}
(J^{\mu\nu}J^{\rho\sigma})^\alpha{}_\beta
=-\eta^{\mu\alpha}\eta^{\rho\nu}\delta^\sigma{}_\beta
+\eta^{\mu\alpha}\eta^{\sigma\nu}\delta^\rho{}_\beta
+\eta^{\nu\alpha}\eta^{\rho\mu}\delta^\sigma{}_\beta
-\eta^{\nu\alpha}\eta^{\sigma\mu}\delta^\rho{}_\beta.
\end{equation*}
四项的物理含义：第一项和第三项含 $\delta^\sigma{}_\beta$（输出指标是 $\sigma$），第二项和第四项含 $\delta^\rho{}_\beta$（输出指标是 $\rho$）。

交换 $(\mu,\nu)\leftrightarrow(\rho,\sigma)$ 得 $J^{\rho\sigma}J^{\mu\nu}$。
\begin{equation*}
(J^{\rho\sigma}J^{\mu\nu})^\alpha{}_\beta
=-\eta^{\rho\alpha}\eta^{\mu\sigma}\delta^\nu{}_\beta
+\eta^{\rho\alpha}\eta^{\nu\sigma}\delta^\mu{}_\beta
+\eta^{\sigma\alpha}\eta^{\mu\rho}\delta^\nu{}_\beta
-\eta^{\sigma\alpha}\eta^{\nu\rho}\delta^\mu{}_\beta.
\end{equation*}

相减并整理。把含 $\delta^\sigma{}_\beta$ 的项合并：
\begin{equation*}
(-\eta^{\mu\alpha}\eta^{\rho\nu}+\eta^{\nu\alpha}\eta^{\rho\mu})\delta^\sigma{}_\beta.
\end{equation*}
利用 $\eta$ 对称性改写为 $\eta^{\rho\mu}$ 的系数：第一项 $-\eta^{\mu\alpha}\eta^{\rho\nu}=-\eta^{\rho\nu}\eta^{\mu\alpha}$，第二项 $\eta^{\nu\alpha}\eta^{\rho\mu}=\eta^{\rho\mu}\eta^{\nu\alpha}$。这恰好可以拆成
\begin{equation*}
\eta^{\mu\rho}\big(\eta^{\nu\alpha}\delta^\sigma{}_\beta\big)-\eta^{\mu\sigma}\big(\eta^{\nu\alpha}\delta^\rho{}_\beta\big)-\eta^{\nu\rho}\big(\eta^{\mu\alpha}\delta^\sigma{}_\beta\big)+\eta^{\nu\sigma}\big(\eta^{\mu\alpha}\delta^\rho{}_\beta\big).
\end{equation*}
每一项的括号内正是某个 $J$ 的分量。例如 $\eta^{\nu\alpha}\delta^\sigma{}_\beta-\eta^{\sigma\alpha}\delta^\nu{}_\beta$ 是 $(J^{\nu\sigma})^\alpha{}_\beta/i$。因此
\begin{equation*}
[J^{\mu\nu},J^{\rho\sigma}]^\alpha{}_\beta
=i\Big(\eta^{\mu\rho}(J^{\nu\sigma})^\alpha{}_\beta-\eta^{\mu\sigma}(J^{\nu\rho})^\alpha{}_\beta-\eta^{\nu\rho}(J^{\mu\sigma})^\alpha{}_\beta+\eta^{\nu\sigma}(J^{\mu\rho})^\alpha{}_\beta\Big).
\end{equation*}
\end{derivation}

由于 Lie 代数的结构常数与具体忠实表示无关（这是 Lie 理论的基本定理：同一个抽象 Lie 代数的任何两个忠实表示给出相同的结构常数），矢量表示中得到的交换关系对一般 Lorentz 生成元 $M^{\mu\nu}$ 同样成立：
\begin{equation}\label{eq:ms-poincare-MM}
\boxed{
[M^{\mu\nu},M^{\rho\sigma}]=i\big(\eta^{\mu\rho}M^{\nu\sigma}-\eta^{\mu\sigma}M^{\nu\rho}-\eta^{\nu\rho}M^{\mu\sigma}+\eta^{\nu\sigma}M^{\mu\rho}\big).
}
\end{equation}
\eqnrefs{eq:ms-poincare-PP}、\eqnrefs{eq:ms-poincare-MP}、\eqnrefs{eq:ms-poincare-MM} 合起来构成 Poincaré Lie 代数。三条关系全部由统一的半直积 Lie 括号\eqnrefs{eq:ms-poincare-unified-bracket} 一次性导出，这正是统一推导的优势：不需要分别猜三组公式的右端，一个群结构给出全部。

物理上，这三条关系有清晰的含义。$[P,P]=0$ 说不同方向的平移互不干扰。$[M,P]\neq 0$ 说 Lorentz 变换会把平移方向旋转到新方向——这就是为什么能量（$P^0$）和动量（$P^i$）在 推进 下会混合。$[M,M]\neq 0$ 说两个旋转的合成一般不是简单相加——例如绕 $x$ 轴旋转再绕 $y$ 轴旋转，与反过来做，结果差一个绕 $z$ 轴的旋转。

\section{Poincaré 群在场上的表示}
到目前为止，$P^\mu$ 和 $M^{\mu\nu}$ 是作用在量子态空间上的抽象算符。但在场论中，基本对象不是态矢量而是场 $\Phi^A(x)$——一组时空依赖的分量，$A$ 标记内部自由度。物理上，场可以是标量（如 Higgs 场 $\phi(x)$）、矢量（如电磁势 $A_\mu(x)$）、旋量（如 Dirac 场 $\psi_\alpha(x)$）等。当坐标系做 Poincaré 变换时，场不仅因为坐标移动而改变（"被动"效应），还因为自身分量在内部空间中旋转而改变（"主动"效应）。后者正是 Lorentz 群的有限维表示 $D(\Lambda)$。
采用被动约定：坐标变换 $x\to x'=\Lambda x+a$ 下，场的变换律为
\begin{equation}\label{eq:ms-poincare-field-transform}
\boxed{
U(\Lambda,a)\Phi^A(x)U^{-1}(\Lambda,a)=D(\Lambda^{-1})^A{}_B\,\Phi^B(\Lambda x+a).
}
\end{equation}
右边有两个因素。$\Phi^B(\Lambda x+a)$ 是坐标移动带来的——场在新坐标处的值。$D(\Lambda^{-1})^A{}_B$ 是内部表示矩阵——场分量本身在 Lorentz 变换下的旋转。注意是 $\Lambda^{-1}$ 而非 $\Lambda$，这是被动变换的约定。

$D$ 必须是 Lorentz 群的表示，即 $D(\Lambda_2\Lambda_1)=D(\Lambda_2)D(\Lambda_1)$。原因很简单：连续做两次变换 $\Lambda_1$ 再 $\Lambda_2$，等价于一次 $\Lambda_2\Lambda_1$，表示矩阵必须按同一规则合成。
写 $\Lambda=I+\omega$，一般有限维表示展开为
\begin{equation}\label{eq:ms-poincare-D-expand}
D(I+\omega)=I-\frac{i}{2}\omega_{\mu\nu}\Sigma^{\mu\nu},
\end{equation}
其中 $\Sigma^{\mu\nu}=-\Sigma^{\nu\mu}$ 是表示空间上的矩阵（不是量子算符），满足与 $M^{\mu\nu}$ 相同的 Lorentz Lie 代数。因子 $-i/2$ 与\eqnrefs{eq:ms-poincare-U-expand} 一致。

定义位移矢量 $\xi^\mu=\epsilon^\mu+\omega^\mu{}_\nu x^\nu$，它把无穷小 Poincaré 变换的坐标移动写成 $x'^\mu=x^\mu+\xi^\mu$。由 Taylor 展开
\begin{equation*}
\Phi^B(\Lambda x+a)=\Phi^B(x)+\xi^\rho\partial_\rho\Phi^B(x)+O(\xi^2),
\end{equation*}
以及 $D(\Lambda^{-1})=D(I-\omega)=I+\frac{i}{2}\omega_{\mu\nu}\Sigma^{\mu\nu}$，代入\eqnrefs{eq:ms-poincare-field-transform} 得到无穷小场变分
\begin{equation}\label{eq:ms-poincare-delta-Phi}
\boxed{
\delta\Phi^A=\xi^\rho\partial_\rho\Phi^A+\frac{i}{2}\omega_{\mu\nu}(\Sigma^{\mu\nu})^A{}_B\Phi^B.
}
\end{equation}
第一项是轨道部分（坐标移动），第二项是内禀部分（场分量旋转）。
纯 Lorentz 变换（$\epsilon=0$）下，$\xi^\mu=\omega^\mu{}_\nu x^\nu$，场变分为
\begin{equation*}
\delta_\omega\Phi=\omega^\mu{}_\nu x^\nu\partial_\mu\Phi+\frac{i}{2}\omega_{\mu\nu}\Sigma^{\mu\nu}\Phi.
\end{equation*}
右端两项分别来自坐标移动和内部旋转。为了把两者统一在一个生成元下，需要把轨道部分也改写成 $\omega_{\mu\nu}$ 的缩并形式。

\begin{derivation}[把轨道部分改写为 Lorentz 生成元的缩并]
轨道部分 $\omega^\mu{}_\nu x^\nu\partial_\mu$ 中，$\omega_{\mu\nu}$ 是反对称参数（$\omega_{\mu\nu}=-\omega_{\nu\mu}$）。目标是把轨道部分写成 $-\frac{i}{2}\omega_{\mu\nu}L^{\mu\nu}$ 的形式，从而与内禀部分 $\frac{i}{2}\omega_{\mu\nu}\Sigma^{\mu\nu}$ 合并为统一的生成元 $M^{\mu\nu}=L^{\mu\nu}+\Sigma^{\mu\nu}$。关键工具是反对称参数与对称张量缩并为零。

降指标。轨道部分中 $\omega^\mu{}_\nu$ 带一个上指标和一个下指标。用 Minkowski 度规降上指标：$\omega^\mu{}_\nu=\eta^{\mu\alpha}\omega_{\alpha\nu}$。代入：
\begin{equation*}
\omega^\mu{}_\nu x^\nu\partial_\mu=\eta^{\mu\alpha}\omega_{\alpha\nu}x^\nu\partial_\mu=\omega_{\alpha\nu}x^\nu\partial^\alpha,
\end{equation*}
最后一步把 $\eta^{\mu\alpha}$ 作用到 $\partial_\mu$ 上得到 $\partial^\alpha\equiv\eta^{\mu\alpha}\partial_\mu$。现在 $\omega_{\alpha\nu}$ 是全下指标，反对称性 $\omega_{\alpha\nu}=-\omega_{\nu\alpha}$ 显式可见。

反对称投影。关键观察：$\omega_{\alpha\nu}$ 反对称，而 $x^\nu\partial^\alpha$ 可以分解为对称和反对称两部分：
\begin{equation*}
x^\nu\partial^\alpha=\frac{1}{2}(x^\nu\partial^\alpha+x^\alpha\partial^\nu)+\frac{1}{2}(x^\nu\partial^\alpha-x^\alpha\partial^\nu).
\end{equation*}
对称部分与反对称 $\omega_{\alpha\nu}$ 缩并为零。验证：
\begin{equation*}
\omega_{\alpha\nu}\cdot\frac{1}{2}(x^\nu\partial^\alpha+x^\alpha\partial^\nu)=\frac{1}{2}(\omega_{\alpha\nu}x^\nu\partial^\alpha+\omega_{\alpha\nu}x^\alpha\partial^\nu).
\end{equation*}
第二项交换哑指标 $\alpha\leftrightarrow\nu$：$\omega_{\nu\alpha}x^\nu\partial^\alpha=-\omega_{\alpha\nu}x^\nu\partial^\alpha$（利用 $\omega_{\nu\alpha}=-\omega_{\alpha\nu}$）。因此两项相加为零。只有反对称部分贡献：
\begin{equation*}
\omega_{\alpha\nu}x^\nu\partial^\alpha=\omega_{\alpha\nu}\cdot\frac{1}{2}(x^\nu\partial^\alpha-x^\alpha\partial^\nu).
\end{equation*}

提出因子并定义轨道角动量。把哑指标 $(\alpha,\nu)$ 换回 $(\mu,\nu)$：
\begin{equation*}
\omega^\mu{}_\nu x^\nu\partial_\mu=\frac{1}{2}\omega_{\mu\nu}(x^\nu\partial^\mu-x^\mu\partial^\nu)=-\frac{1}{2}\omega_{\mu\nu}(x^\mu\partial^\nu-x^\nu\partial^\mu).
\end{equation*}
为了与量子力学中 Hermite 生成元的约定一致（无穷小算符 $U=I-\frac{i}{2}\omega_{\mu\nu}M^{\mu\nu}$ 中 $M^{\mu\nu}$ 应为 Hermite），引入因子 $i$：
\begin{equation*}
-\frac{1}{2}\omega_{\mu\nu}(x^\mu\partial^\nu-x^\nu\partial^\mu)=-\frac{i}{2}\omega_{\mu\nu}\cdot i(x^\mu\partial^\nu-x^\nu\partial^\mu).
\end{equation*}
定义轨道角动量算符
\begin{equation}\label{eq:ms-poincare-L-def}
L^{\mu\nu}=i(x^\mu\partial^\nu-x^\nu\partial^\mu).
\end{equation}
于是轨道部分 = $-\frac{i}{2}\omega_{\mu\nu}L^{\mu\nu}$，内禀部分 = $\frac{i}{2}\omega_{\mu\nu}\Sigma^{\mu\nu}$，两者形式统一。

还需验证 $L^{\mu\nu}$ 的 Hermite 性。$x^\mu$ 和 $\partial^\nu$ 在量子力学中是共轭算符（$(x^\mu)^\dagger=x^\mu$，$(\partial^\nu)^\dagger=-\partial^\nu$ 在适当内积下）。因此
\begin{equation*}
(L^{\mu\nu})^\dagger=-i(x^\nu\partial^\mu-x^\mu\partial^\nu)=i(x^\mu\partial^\nu-x^\nu\partial^\mu)=L^{\mu\nu},
\end{equation*}
确认 $L^{\mu\nu}$ 是 Hermite 算符，对应物理可观测量（角动量）。
\end{derivation}

合并后，场上的总 Lorentz 生成元自然分为轨道与内禀两部分：
\begin{equation*}\boxed{M^{\mu\nu}_{\text{field}}=L^{\mu\nu}+\Sigma^{\mu\nu}.}
\end{equation*}
这个分解的物理意义非常重要。$L^{\mu\nu}$ 只涉及坐标 $x$ 和导数 $\partial$，与场的内部类型无关——无论是标量、矢量还是旋量，轨道部分都一样。$\Sigma^{\mu\nu}$ 只作用在场的内部指标上，与时空位置无关——标量的 $\Sigma=0$，矢量的 $\Sigma$ 是 $d\times d$ 矩阵，旋量的 $\Sigma$ 是 Clifford 代数构造。不同类型的场，只是选择了不同的 $\Sigma^{\mu\nu}$。

\section{任意阶张量与旋量的 Lorentz 生成元}

上一节把场的 Lorentz 生成元分解为 $M=L+\Sigma$，其中 $\Sigma$ 由场的类型决定。本节逐一计算常见场类型的 $\Sigma^{\mu\nu}$。
标量场是最简单的情形：$\phi'(x')=\phi(x)$，场分量不旋转，$D(\Lambda)=1$。因此
\begin{equation*}
\boxed{\Sigma^{\mu\nu}=0\quad\text{(标量)}.}
\end{equation*}
标量场的 Lorentz 生成元纯轨道：$M^{\mu\nu}=L^{\mu\nu}=i(x^\mu\partial^\nu-x^\nu\partial^\mu)$。

四矢量场 $V^\rho$ 在 Lorentz 变换下按 $V'^\rho=\Lambda^\rho{}_\sigma V^\sigma$ 变换。无穷小形式 $\delta V^\rho=\omega^\rho{}_\sigma V^\sigma$ 应与\eqnrefs{eq:ms-poincare-D-expand} 给出的 $\delta V^\rho=-\frac{i}{2}\omega_{\mu\nu}(\Sigma^{\mu\nu})^\rho{}_\sigma V^\sigma$ 一致。比较 $\omega_{\mu\nu}$ 的系数（注意 $\omega^\rho{}_\sigma=\eta^{\rho\mu}\omega_{\mu\sigma}$）得到
\begin{equation*}\boxed{
(\Sigma^{\mu\nu})^\rho{}_\sigma=i(\eta^{\mu\rho}\delta^\nu{}_\sigma-\eta^{\nu\rho}\delta^\mu{}_\sigma)\quad\text{(矢量)}.
}
\end{equation*}
这与矢量表示的 Lorentz 生成元\eqnrefs{eq:ms-poincare-JV} 一致——这正是预期的，因为矢量场按矢量表示变换。
一般 $(r,s)$ 张量 $T^{\alpha_1\cdots\alpha_r}{}_{\beta_1\cdots\beta_s}$ 有 $r$ 个上指标和 $s$ 个下指标。Lorentz 变换对每一个指标独立作用：每个上指标贡献一个矢量表示的 $\Sigma$，每个下指标贡献一个协变矢量表示的 $\Sigma$。总生成元是各指标生成元之和。显式地：
\begin{equation*}\boxed{
\begin{aligned}
(\Sigma^{\mu\nu}T)^{\alpha_1\cdots\alpha_r}{}_{\beta_1\cdots\beta_s}
={}&i\sum_{k=1}^{r}\Big(\eta^{\mu\alpha_k}T^{\cdots\nu\cdots}{}_{\cdots}-\eta^{\nu\alpha_k}T^{\cdots\mu\cdots}{}_{\cdots}\Big)\\
&+i\sum_{\ell=1}^{s}\Big(\delta^\mu{}_{\beta_\ell}T^{\cdots}{}_{\cdots\nu\cdots}-\delta^\nu{}_{\beta_\ell}T^{\cdots}{}_{\cdots\mu\cdots}\Big),
\end{aligned}
}
\end{equation*}
其中省略号表示未被替换的指标保持不变。例如对二阶逆变张量 $T^{\rho\sigma}$：
\begin{equation*}
(\Sigma^{\mu\nu}T)^{\rho\sigma}=i\big(\eta^{\mu\rho}T^{\nu\sigma}-\eta^{\nu\rho}T^{\mu\sigma}+\eta^{\mu\sigma}T^{\rho\nu}-\eta^{\nu\sigma}T^{\rho\mu}\big).
\end{equation*}
这就是 $\Sigma^{\mu\nu}_{(2)}=\Sigma^{\mu\nu}_{(1)}\otimes I+I\otimes\Sigma^{\mu\nu}_{(1)}$——张量积表示的生成元是各因子生成元之和。
旋量场的 $\Sigma^{\mu\nu}$ 不能从普通张量指标构造，因为旋量不属于任何 $SO(1,d-1)$ 的张量表示。它属于 Lorentz 群的双覆盖 $Spin(1,d-1)$ 的旋量表示。构造旋量表示的出发点是 Clifford 代数
\begin{equation*}\boxed{\{\gamma^\mu,\gamma^\nu\}=2\eta^{\mu\nu},}
\end{equation*}
它要求 $\gamma^\mu$ 是 $2^{\lfloor d/2\rfloor}\times 2^{\lfloor d/2\rfloor}$ 矩阵（$d=4$ 时是 $4\times4$）。旋量表示 $D(\Lambda)$ 必须保持 Clifford 代数协变：
\begin{equation}\label{eq:ms-poincare-clifford-cov}
D(\Lambda)\gamma^\rho D^{-1}(\Lambda)=\Lambda^\rho{}_\sigma\gamma^\sigma.
\end{equation}
这个条件的物理意义是：$\gamma^\mu$ 作为"矢量矩阵"在 Lorentz 变换下应按矢量表示旋转，而 $D(\Lambda)$ 是实现这一旋转的旋量表示矩阵。

\begin{derivation}[从 Clifford 代数导出旋量表示的 Lorentz 生成元]
把 $D(\Lambda)=I-\frac{i}{2}\omega_{\mu\nu}\Sigma^{\mu\nu}$ 代入协变条件\eqnrefs{eq:ms-poincare-clifford-cov}，比较 $\omega$ 的一阶项，解出 $\Sigma^{\mu\nu}$。

展开左边。把 $D$ 和 $D^{-1}$ 代入\eqnrefs{eq:ms-poincare-clifford-cov} 左边：
\begin{equation*}
(I-\tfrac{i}{2}\omega_{\mu\nu}\Sigma^{\mu\nu})\gamma^\rho(I+\tfrac{i}{2}\omega_{\alpha\beta}\Sigma^{\alpha\beta})
=\gamma^\rho-\frac{i}{2}\omega_{\mu\nu}\Sigma^{\mu\nu}\gamma^\rho+\frac{i}{2}\omega_{\alpha\beta}\gamma^\rho\Sigma^{\alpha\beta}+O(\omega^2).
\end{equation*}
合并后两项为交换子：
\begin{equation*}
=\gamma^\rho-\frac{i}{2}\omega_{\mu\nu}[\Sigma^{\mu\nu},\gamma^\rho]+O(\omega^2).
\end{equation*}

展开右边。$\Lambda^\rho{}_\sigma=\delta^\rho{}_\sigma+\omega^\rho{}_\sigma$，故
\begin{equation*}
\Lambda^\rho{}_\sigma\gamma^\sigma=\gamma^\rho+\omega^\rho{}_\sigma\gamma^\sigma.
\end{equation*}
把 $\omega^\rho{}_\sigma=\eta^{\rho\mu}\omega_{\mu\sigma}$ 代入，并利用 $\omega_{\mu\sigma}$ 反对称性把单个 $\eta^{\rho\mu}\omega_{\mu\sigma}\gamma^\sigma$ 改写为反对称缩并：
\begin{equation*}
\omega^\rho{}_\sigma\gamma^\sigma=\eta^{\rho\mu}\omega_{\mu\nu}\gamma^\nu=\frac{1}{2}\omega_{\mu\nu}(\eta^{\rho\mu}\gamma^\nu-\eta^{\rho\nu}\gamma^\mu)=\omega_{\mu\nu}\eta^{\mu\rho}\gamma^\nu.
\end{equation*}
最后一步用了 $\omega_{\mu\nu}$ 反对称：$\frac{1}{2}(\eta^{\rho\mu}\gamma^\nu-\eta^{\rho\nu}\gamma^\mu)=\eta^{\rho\mu}\gamma^\nu$（交换 $\mu\leftrightarrow\nu$ 在第二项中利用 $\omega_{\nu\mu}=-\omega_{\mu\nu}$ 恢复）。

比较 $\omega_{\mu\nu}$ 的系数，得到
\begin{equation}\label{eq:ms-poincare-sigma-gamma-comm}
-\frac{i}{2}[\Sigma^{\mu\nu},\gamma^\rho]=\eta^{\mu\rho}\gamma^\nu-\eta^{\nu\rho}\gamma^\mu.
\end{equation}

为验证 $\Sigma^{\mu\nu}=\frac{i}{4}[\gamma^\mu,\gamma^\nu]$ 满足此式，先由 Clifford 代数 $\{\gamma^\mu,\gamma^\nu\}=2\eta^{\mu\nu}$ 得到辅助恒等式
\begin{align*}
[\gamma^\mu\gamma^\nu,\gamma^\rho]
&=\gamma^\mu\gamma^\nu\gamma^\rho-\gamma^\rho\gamma^\mu\gamma^\nu\\
&=\gamma^\mu\{\gamma^\nu,\gamma^\rho\}-\{\gamma^\mu,\gamma^\rho\}\gamma^\nu\\
&=2\eta^{\nu\rho}\gamma^\mu-2\eta^{\mu\rho}\gamma^\nu.
\end{align*}
第一个等号把 $\gamma^\nu\gamma^\rho$ 拆成 $\{\gamma^\nu,\gamma^\rho\}-\gamma^\rho\gamma^\nu$，第二个等号用 Clifford 代数。同理
\begin{equation*}
[\gamma^\nu\gamma^\mu,\gamma^\rho]=2\eta^{\mu\rho}\gamma^\nu-2\eta^{\nu\rho}\gamma^\mu.
\end{equation*}
相减得
\begin{equation*}
[[\gamma^\mu,\gamma^\nu],\gamma^\rho]=4(\eta^{\nu\rho}\gamma^\mu-\eta^{\mu\rho}\gamma^\nu).
\end{equation*}
因此
\begin{equation*}
[\Sigma^{\mu\nu},\gamma^\rho]=\frac{i}{4}[[\gamma^\mu,\gamma^\nu],\gamma^\rho]=\frac{i}{4}\cdot 4(\eta^{\nu\rho}\gamma^\mu-\eta^{\mu\rho}\gamma^\nu)=i(\eta^{\nu\rho}\gamma^\mu-\eta^{\mu\rho}\gamma^\nu).
\end{equation*}
代入\eqnrefs{eq:ms-poincare-sigma-gamma-comm} 左边：
\begin{equation*}
-\frac{i}{2}\cdot i(\eta^{\nu\rho}\gamma^\mu-\eta^{\mu\rho}\gamma^\nu)=\frac{1}{2}(\eta^{\nu\rho}\gamma^\mu-\eta^{\mu\rho}\gamma^\nu)=\eta^{\mu\rho}\gamma^\nu-\eta^{\nu\rho}\gamma^\mu,
\end{equation*}
与右边一致。
\end{derivation}

因此 Dirac 旋量的 Lorentz 生成元
\begin{equation}\label{eq:ms-poincare-Sigma-dirac}
\boxed{\Sigma^{\mu\nu}=\frac{i}{4}[\gamma^\mu,\gamma^\nu]\quad\text{(Dirac 旋量)}.}
\end{equation}
Dirac 场上的总 Lorentz 生成元
\begin{equation*}\boxed{
M^{\mu\nu}=i(x^\mu\partial^\nu-x^\nu\partial^\mu)+\frac{i}{4}[\gamma^\mu,\gamma^\nu].
}
\end{equation*}
第一项是轨道角动量，第二项是自旋角动量。对 $d=4$，自旋部分可以分解为 $\frac{i}{4}[\gamma^\mu,\gamma^\nu]=\frac{i}{2}\sigma^{\mu\nu}$，其中 $\sigma^{\mu\nu}=\frac{1}{2}[\gamma^\mu,\gamma^\nu]$ 是标准的 Dirac 自旋矩阵。
在 $d=4$ 时，Lorentz 群的双覆盖 $Spin(1,3)\simeq SL(2,\mathbb{C})$。有限维不可约表示由一对半整数 $(j_L,j_R)$ 标记，其中 $j_L,j_R=0,\frac{1}{2},1,\ldots$。这个标记的来源是：$SL(2,\mathbb{C})$ 同构于 $SU(2)_L\times SU(2)_R$（在复化意义上），每个因子有自己的角动量量子数。常见场的表示对应关系：

\begin{center}
\begin{tabular}{ll}
\hline
场 & Lorentz 表示 \\
\hline
标量 & $(0,0)$ \\
左手 Weyl & $(\frac{1}{2},0)$ \\
右手 Weyl & $(0,\frac{1}{2})$ \\
Dirac & $(\frac{1}{2},0)\oplus(0,\frac{1}{2})$ \\
四矢量 & $(\frac{1}{2},\frac{1}{2})$ \\
自对偶二形式 & $(1,0)$ \\
反自对偶二形式 & $(0,1)$ \\
对称无迹 秩-2 & $(1,1)$ \\
\hline
\end{tabular}
\end{center}

这张表是后续自旋分类的基础。注意 Dirac 旋量是可约表示——它由左手和右手 Weyl 旋量直和而成。质量为零时这两个分量解耦，各自独立描述不同螺旋度的粒子；质量不为零时质量项把它们耦合，这就是下一章要详细讨论的 手征性-螺旋度 关系。

\section{Wigner 分类：场表示与粒子自旋}

到这里有一个必须强调的区别：场的 Lorentz 表示不等于单粒子的 Poincaré 不可约表示。场 $\Phi^A(x)$ 的指标 $A$ 属于有限维 Lorentz 表示，它描述的是场在 Lorentz 变换下的变换律。但物理单粒子态是 Hilbert 空间中的向量，它们按照 Poincaré 群的\emph{无穷维}酉表示分类——Wigner 1939 年的分类方案。同一个场可以产生多种粒子态（例如矢量场 $A_\mu$ 在无质量时只产生 $h=\pm1$ 两个螺旋度的光子，而非 4 个分量对应 4 种粒子）。

Wigner 分类的关键是两个 Casimir 不变量——它们与所有 Poincaré 生成元对易，因此在不可约表示上是常数。第一 Casimir 是
\begin{equation*}\boxed{C_1=P^2=P_\mu P^\mu,}
\end{equation*}
对质量为 $m$ 的粒子给出 $P^2=-m^2$（本文"多为正"度规下）。第二 Casimir 由 Pauli--Lubanski 向量
\begin{equation*}W^\mu=\frac{1}{2}\epsilon^{\mu\nu\rho\sigma}P_\nu M_{\rho\sigma}
\end{equation*}
构造：
\begin{equation*}\boxed{C_2=W^2=W_\mu W^\mu.}
\end{equation*}
对 有质量 自旋-$s$ 粒子，$W^2=m^2 s(s+1)$。无质量粒子则由 小群 分类——有质量 的 小群 是 $SO(3)$，无质量 的是 $ISO(2)$，两者的表示结构截然不同。这些细节留到下一章系统讨论。

\section{最小作用量原理与 Euler--Lagrange 方程}

有了场的变换律，下一步问：什么样的场是物理允许的？经典场论的回答是作用量驻定。考虑一阶局域场论
\begin{equation}
S[\Phi]=\int \dd^dx\,\mathcal{L}(\Phi^A,\partial_\mu\Phi^A),
\end{equation}
其中拉氏密度 $\mathcal{L}$ 依赖于场和它的一阶导数。物理场是使 $\delta S=0$（对边界消失的任意变分）的场。

\begin{derivation}[Euler--Lagrange 方程的推导]
对作用量变分，分部积分把 $\delta\Phi$ 的导数移走，使变分只以 $\delta\Phi^A$（而非 $\partial_\mu\delta\Phi^A$）的形式出现。然后利用 $\delta\Phi^A$ 的任意性读出方程。

写出变分。场的任意变分 $\Phi^A\to\Phi^A+\delta\Phi^A$ 引起
\begin{equation*}
\delta S=\int \dd^dx\left[\frac{\partial\mathcal{L}}{\partial\Phi^A}\delta\Phi^A+\frac{\partial\mathcal{L}}{\partial(\partial_\mu\Phi^A)}\partial_\mu\delta\Phi^A\right].
\end{equation*}
第一项来自 $\mathcal{L}$ 对场的显式依赖，第二项来自 $\mathcal{L}$ 对场导数的依赖。

定义共轭动量。为简洁起见定义
\begin{equation*}
\Pi_A^\mu\equiv\frac{\partial\mathcal{L}}{\partial(\partial_\mu\Phi^A)}.
\end{equation*}
物理上 $\Pi_A^0$ 是场的正则动量，$\Pi_A^i$ 是流的空间分量。

分部积分。第二项中 $\partial_\mu\delta\Phi^A$ 是变分的导数。用 Leibniz 法则把它移到 $\Pi_A^\mu$ 上：
\begin{equation*}
\Pi_A^\mu\partial_\mu\delta\Phi^A=\partial_\mu(\Pi_A^\mu\delta\Phi^A)-(\partial_\mu\Pi_A^\mu)\delta\Phi^A.
\end{equation*}
第一项是全散度，积分后化为边界项。

代入并丢掉边界项。
\begin{equation*}
\delta S=\int \dd^dx\left[\left(\frac{\partial\mathcal{L}}{\partial\Phi^A}-\partial_\mu\Pi_A^\mu\right)\delta\Phi^A+\partial_\mu(\Pi_A^\mu\delta\Phi^A)\right].
\end{equation*}
若变分 $\delta\Phi^A$ 在边界上消失（或场在无穷远衰减足够快），散度项不贡献。

利用 $\delta\Phi^A$ 任意性。剩余积分中 $\delta\Phi^A$ 任意，要使 $\delta S=0$ 对所有 $\delta\Phi^A$ 成立，系数必须为零：
\begin{equation}\label{eq:ms-poincare-EL}
\boxed{
\frac{\partial\mathcal{L}}{\partial\Phi^A}-\partial_\mu\frac{\partial\mathcal{L}}{\partial(\partial_\mu\Phi^A)}=0.
}
\end{equation}
这就是 Euler--Lagrange 方程。对每个场分量 $A$ 有一个方程。

推广与边界：高阶导数 Lagrangian 给出 Ostrogradsky 方程；退化 Lagrangian 的约束由 Dirac--Bergmann 算法处理，规范对称性来自第一类约束。变分边界项 $\theta^\mu=\frac{\partial\mathcal{L}}{\partial(\partial_\mu\Phi^A)}\delta\Phi^A$ 即 Noether 流的变分形式，给出荷的定义 $Q=\int_\Sigma\theta^0\,\dd^3x$；Euler--Lagrange 方程与 Noether 定理同源。
\end{derivation}

\section{Noether 第一定理的一般形式}

Noether 定理把连续对称性与守恒律联系起来。核心思想很简单：如果某种连续变换使作用量不变（最多差一个边界项），那么对应的变分恒等式自动给出一个散度为零的流——守恒流。

首先把任意变分写成标准形式。利用 Euler--Lagrange 方程的推导过程，任何变分都可以写成
\begin{equation}\label{eq:ms-poincare-delta-L}
\boxed{
\delta\mathcal{L}=E_A\,\delta\Phi^A+\partial_\mu\Theta^\mu,
}
\end{equation}
其中 $E_A=\frac{\partial\mathcal{L}}{\partial\Phi^A}-\partial_\mu\Pi_A^\mu$ 是 EL 表达式（$E_A=0$ 即运动方程），$\Theta^\mu=\Pi_A^\mu\delta\Phi^A$ 是 Symplectic 势。这个分解的物理意义：变分的效果分为两部分——一部分正比于运动方程（在壳 为零），另一部分是全散度（边界项）。

现在假设某种连续变换使拉氏量的变分是一个全散度：
\begin{equation}
\delta\mathcal{L}=\partial_\mu K^\mu.
\end{equation}
这种变换称为"准对称"（quasi-对称）。例如平移变换使 $\delta\mathcal{L}=\epsilon^\nu\partial_\nu\mathcal{L}=\partial_\mu(\epsilon^\mu\mathcal{L})$，故 $K^\mu=\epsilon^\mu\mathcal{L}$。把\eqnrefs{eq:ms-poincare-delta-L} 与准对称条件结合：
\begin{equation*}
E_A\,\delta\Phi^A+\partial_\mu\Theta^\mu=\partial_\mu K^\mu,
\end{equation*}
即
\begin{equation*}
E_A\,\delta\Phi^A+\partial_\mu(\Theta^\mu-K^\mu)=0.
\end{equation*}
定义 Noether 流
\begin{equation}\label{eq:ms-poincare-J-def}
\boxed{J^\mu=\Theta^\mu-K^\mu=\Pi_A^\mu\delta\Phi^A-K^\mu,}
\end{equation}
则 off-shell 恒等式为
\begin{equation*}\boxed{\partial_\mu J^\mu=-E_A\,\delta\Phi^A.}
\end{equation*}
在运动方程成立时（在壳，$E_A=0$），右端为零：
\begin{equation}\label{eq:ms-poincare-noether}
\boxed{\partial_\mu J^\mu=0\quad\text{(on-shell)}.}
\end{equation}
这就是 Noether 第一定理：每一个连续准对称性给出一个 在壳 守恒流。

\section{Poincaré 对称性与能动量守恒}

现在把 Noether 定理应用到 Poincaré 对称性的两个子群：平移和 Lorentz。
纯平移 $x^\mu\to x^\mu+\epsilon^\mu$ 给出 $\delta\Phi^A=\epsilon^\nu\partial_\nu\Phi^A$。若拉氏量不显含 $x$（时空均匀），则
\begin{equation*}
\delta\mathcal{L}=\epsilon^\nu\partial_\nu\mathcal{L}=\partial_\mu(\epsilon^\mu\mathcal{L}),
\end{equation*}
故 $K^\mu=\epsilon^\mu\mathcal{L}$。

\begin{derivation}[能动量张量的推导]
把平移变分 $\delta\Phi^A=\epsilon^\nu\partial_\nu\Phi^A$ 代入 Noether 流\eqnrefs{eq:ms-poincare-J-def}，提出任意参数 $\epsilon^\nu$，剩余部分定义为能动量张量。关键步骤是把 $K^\mu=\epsilon^\mu\mathcal{L}$ 中的指标 $\mu$ 改写为 $\epsilon^\nu$ 的缩并形式，使所有项共享同一个任意参数。

Symplectic 势。平移变分 $\delta\Phi^A=\epsilon^\nu\partial_\nu\Phi^A$ 代入 $\Theta^\mu=\Pi_A^\mu\delta\Phi^A$：
\begin{equation*}
\Theta^\mu=\Pi_A^\mu\,\epsilon^\nu\partial_\nu\Phi^A=\epsilon^\nu\Pi_A^\mu\partial_\nu\Phi^A.
\end{equation*}
这里 $\epsilon^\nu$ 是任意常矢量（平移参数），可以提到缩并外面。$\Pi_A^\mu\partial_\nu\Phi^A$ 是"动量密度乘速度"的结构——$\Pi_A^\mu$ 是共轭动量（场的 $\mu$ 方向流），$\partial_\nu\Phi^A$ 是场在 $\nu$ 方向的变化率。

Noether 流。由\eqnrefs{eq:ms-poincare-J-def}，$J^\mu=\Theta^\mu-K^\mu$。平移变换使拉氏量变分为全散度 $\delta\mathcal{L}=\epsilon^\nu\partial_\nu\mathcal{L}=\partial_\mu(\epsilon^\mu\mathcal{L})$（利用 $\epsilon^\nu$ 常矢量），故 $K^\mu=\epsilon^\mu\mathcal{L}$。Noether 流为
\begin{equation*}
J^\mu=\epsilon^\nu\Pi_A^\mu\partial_\nu\Phi^A-\epsilon^\mu\mathcal{L}.
\end{equation*}
第二项 $\epsilon^\mu\mathcal{L}$ 的自由指标是 $\mu$，而第一项的任意参数指标是 $\nu$。为了合并两项，需要统一任意参数。利用 $\epsilon^\mu=\delta^\mu{}_\nu\epsilon^\nu$（Kronecker δ 转指标）：
\begin{equation*}
\epsilon^\mu\mathcal{L}=\delta^\mu{}_\nu\epsilon^\nu\mathcal{L}=\epsilon^\nu\delta^\mu{}_\nu\mathcal{L}.
\end{equation*}
代入后 Noether 流统一为
\begin{equation*}
J^\mu=\epsilon^\nu\big[\Pi_A^\mu\partial_\nu\Phi^A-\delta^\mu{}_\nu\mathcal{L}\big].
\end{equation*}

定义 正则 能动量张量。因为 $\epsilon^\nu$ 任意，方括号中的量对每个 $\nu$ 独立守恒。定义 正则 能动量张量
\begin{equation}\label{eq:ms-poincare-T-can}
\boxed{
T^\mu{}_\nu=\Pi_A^\mu\partial_\nu\Phi^A-\delta^\mu{}_\nu\mathcal{L}.
}
\end{equation}
于是 $J^\mu_{\text{trans}}=\epsilon^\nu T^\mu{}_\nu$。第一项 $\Pi_A^\mu\partial_\nu\Phi^A$ 是"动量流乘场速度"——物理上是能量-动量流的密度。第二项 $-\delta^\mu{}_\nu\mathcal{L}$ 是拉氏密度的对角项——它来自平移变分使拉氏量本身移动的效应。

守恒律。由 Noether 定理\eqnrefs{eq:ms-poincare-noether}，$\partial_\mu J^\mu=0$ 在壳。代入 $J^\mu=\epsilon^\nu T^\mu{}_\nu$：
\begin{equation*}
0=\partial_\mu J^\mu=\partial_\mu(\epsilon^\nu T^\mu{}_\nu)=\epsilon^\nu\,\partial_\mu T^\mu{}_\nu.
\end{equation*}
最后一步利用 $\epsilon^\nu$ 常矢量（$\partial_\mu\epsilon^\nu=0$）。由于 $\epsilon^\nu$ 任意，每个分量独立为零：
\begin{equation}\label{eq:ms-poincare-T-conservation}
\boxed{\partial_\mu T^\mu{}_\nu=0.}
\end{equation}

一般 $T^{\mu\nu}$ 不对称，但可加全导数项（Belinfante 修正）使 $\Theta^{\mu\nu}$ 对称且守恒，详见下框。
能动量张量也可由作用量对度量的变分定义 $T_{\mu\nu}=\frac{-2}{\sqrt{-g}}\frac{\delta S}{\delta g^{\mu\nu}}$，自动给出对称张量，并在弯曲时空推广为协变守恒 $\nabla_\mu T^{\mu\nu}=0$。
\end{derivation}

物理上，$T^{0}{}_\nu$ 是能量密度（$\nu=0$）和动量密度（$\nu=i$）。$T^{i}{}_\nu$ 是对应的流——能量流和动量流（应力）。守恒方程 $\partial_\mu T^\mu{}_\nu=0$ 说：每个时间步中能量/动量密度的变化，等于通过边界的能量/动量流的负值。
由连续性方程 $\partial_0 T^{0\nu}+\partial_i T^{i\nu}=0$ 对空间积分：
\begin{equation*}
\frac{\dd}{\dd t}\int \dd^{d-1}x\,T^{0\nu}=-\int \dd^{d-1}x\,\partial_i T^{i\nu}.
\end{equation*}
若场在空间无穷远衰减足够快，右面积分为零（Gauss 定理），故
\begin{equation*}\boxed{P^\nu=\int \dd^{d-1}x\,T^{0\nu}}
\end{equation*}
守恒。这就是四动量算符——能量和动量的积分形式。

\section{Lorentz 对称性与角动量守恒}

纯 Lorentz 变换（$\epsilon^\mu=0$）下，$\xi^\mu=\omega^\mu{}_\nu x^\nu$，场变分由\eqnrefs{eq:ms-poincare-delta-Phi} 给出。因为 $\mathcal{L}$ 是 Lorentz 标量，
\begin{equation*}
\delta\mathcal{L}=\xi^\rho\partial_\rho\mathcal{L}.
\end{equation*}
又因 $\partial_\rho\xi^\rho=\omega^\rho{}_\rho=0$（$\omega$ 无迹），故
\begin{equation*}
\delta\mathcal{L}=\partial_\rho(\xi^\rho\mathcal{L}),
\end{equation*}
即 $K^\lambda=\xi^\lambda\mathcal{L}$。

\begin{derivation}[Lorentz Noether 流与角动量守恒]
把 Lorentz 变分代入 Noether 流，分离出轨道部分（含 $x$ 和 $T$）和自旋部分（含 $\Sigma$），分别定义轨道角动量流和自旋流。

Noether 流。
\begin{equation*}
J^\lambda=\Pi_A^\lambda\delta\Phi^A-\xi^\lambda\mathcal{L}
=\omega^\rho{}_\sigma x^\sigma\big(\Pi_A^\lambda\partial_\rho\Phi^A-\delta^\lambda{}_\rho\mathcal{L}\big)+\frac{i}{2}\omega_{\mu\nu}\Pi_A^\lambda(\Sigma^{\mu\nu})^A{}_B\Phi^B.
\end{equation*}

识别能动量张量。第一括号正是 $T^\lambda{}_\rho$（见\eqnrefs{eq:ms-poincare-T-can}）。

轨道部分改写。利用 $\omega_{\mu\nu}$ 反对称性，把 $\omega^\rho{}_\sigma x^\sigma T^\lambda{}_\rho$ 改写为 $\omega_{\mu\nu}$ 的缩并：
\begin{equation*}
\omega^\rho{}_\sigma x^\sigma T^\lambda{}_\rho=-\frac{1}{2}\omega_{\mu\nu}(x^\mu T^{\lambda\nu}-x^\nu T^{\lambda\mu}).
\end{equation*}
这与\eqnrefs{eq:ms-poincare-L-def} 中轨道部分的改写完全平行。

定义轨道与自旋流。
\begin{equation}\label{eq:ms-poincare-L-current}
L^{\lambda\mu\nu}=x^\mu T^{\lambda\nu}-x^\nu T^{\lambda\mu},
\end{equation}
\begin{equation}\label{eq:ms-poincare-S-current}
S^{\lambda\mu\nu}=-i\,\Pi_A^\lambda(\Sigma^{\mu\nu})^A{}_B\Phi^B.
\end{equation}
轨道流 $L^{\lambda\mu\nu}$ 只依赖坐标和能动量张量，对所有场类型一样。自旋流 $S^{\lambda\mu\nu}$ 依赖场的内禀表示 $\Sigma^{\mu\nu}$，标量场为零。

合并。
\begin{equation*}
J^\lambda_{\text{Lorentz}}=-\frac{1}{2}\omega_{\mu\nu}(L^{\lambda\mu\nu}+S^{\lambda\mu\nu}).
\end{equation*}
定义总角动量流
\begin{equation}\label{eq:ms-poincare-J-total}
J^{\lambda\mu\nu}=L^{\lambda\mu\nu}+S^{\lambda\mu\nu},
\end{equation}
由 $\omega_{\mu\nu}$ 任意得
\begin{equation}\label{eq:ms-poincare-J-conservation}
\boxed{\partial_\lambda J^{\lambda\mu\nu}=0.}
\end{equation}

$J^{\lambda\mu\nu}=x^\mu T^{\lambda\nu}-x^\nu T^{\lambda\mu}+S^{\lambda\mu\nu}$，前两项是轨道角动量流（从 $T^{\lambda\nu}$ 构造），$S^{\lambda\mu\nu}=\frac{\partial\mathcal{L}}{\partial(\partial_\lambda\Phi^A)}(\Sigma^{\mu\nu})^A{}_B\Phi^B$ 是自旋流（来自场的内禀 Lorentz 变换 $\delta\Phi^A=\frac12\omega_{\mu\nu}(\Sigma^{\mu\nu})^A{}_B\Phi^B$）。对标量场 $S^{\lambda\mu\nu}=0$（无自旋），对 Dirac 场 $S^{\lambda\mu\nu}=\bar\psi\gamma^\lambda\Sigma^{\mu\nu}\psi$（$\Sigma^{\mu\nu}=\frac14[\gamma^\mu,\gamma^\nu]$，自旋 $1/2$），对矢量场 $S^{\lambda\mu\nu}=F^{\lambda\mu}A^\nu-F^{\lambda\nu}A^\mu$（自旋 $1$）。守恒律 $\partial_\lambda J^{\lambda\mu\nu}=0$ 展开为 $T^{\mu\nu}-T^{\nu\mu}+\partial_\lambda S^{\lambda\mu\nu}=0$，给出能动量张量对称性与自旋流的关系：若 $T^{\mu\nu}$ 已对称（Belinfante 修正后），则 $\partial_\lambda S^{\lambda\mu\nu}=0$（自旋流单独守恒）。

总角动量 $J_i=\frac12\epsilon_{ijk}\int J^{0jk}\,\dd^3x$ 满足 $[J_i,J_j]=i\epsilon_{ijk}J_k$；Poincar\'e 群的两个 Casimir $P^2$（质量）与 $W^2$（自旋，$W_\mu=\frac12\epsilon_{\mu\nu\rho\sigma}J^{\nu\rho}P^\sigma$ 为 Pauli--Lubanski 矢量）完全标记单粒子态，是 Wigner 分类的基础。
\end{derivation}

角动量流由此分解：轨道流 $L^{\lambda\mu\nu}=x^\mu T^{\lambda\nu}-x^\nu T^{\lambda\mu}$（\eqnrefs{eq:ms-poincare-L-current}）、自旋流 $S^{\lambda\mu\nu}=-i\Pi_A^\lambda(\Sigma^{\mu\nu})^A{}_B\Phi^B$（\eqnrefs{eq:ms-poincare-S-current}），总角动量流 $J^{\lambda\mu\nu}=L^{\lambda\mu\nu}+S^{\lambda\mu\nu}$ 满足 $\partial_\lambda J^{\lambda\mu\nu}=0$（\eqnrefs{eq:ms-poincare-J-total}、\eqnrefs{eq:ms-poincare-J-conservation}）。

守恒荷
\begin{equation*}\boxed{
M^{\mu\nu}=\int \dd^{d-1}x\,\big[x^\mu T^{0\nu}-x^\nu T^{0\mu}+S^{0\mu\nu}\big].
}
\end{equation*}
前两项是轨道角动量 $\mathbf{x}\times\mathbf{p}$ 的相对论推广，第三项是自旋角动量。对标量场 $S=0$，总角动量纯轨道。对旋量场，自旋部分贡献 $\frac{1}{2}$ 的内禀自旋。
总角动量流守恒 $\partial_\lambda J^{\lambda\mu\nu}=0$ 展开后给出 正则 能动量张量的一个重要性质。展开：
\begin{equation*}
0=\partial_\lambda(x^\mu T^{\lambda\nu}-x^\nu T^{\lambda\mu}+S^{\lambda\mu\nu}).
\end{equation*}
利用 $\partial_\lambda x^\mu=\delta^\mu{}_\lambda$ 和 $\partial_\lambda T^{\lambda\nu}=0$（能动量守恒）：
\begin{equation*}
0=T^{\mu\nu}-T^{\nu\mu}+\partial_\lambda S^{\lambda\mu\nu}.
\end{equation*}
即
\begin{equation}\label{eq:ms-poincare-T-asym}
\boxed{
T^{\mu\nu}-T^{\nu\mu}=-\partial_\lambda S^{\lambda\mu\nu}.
}
\end{equation}
正则 能动量张量的反对称部分由自旋流的散度决定。对标量场 $S=0$，$T^{\mu\nu}$ 自动对称。对旋量场，$T^{\mu\nu}$ 一般不对称——这正是 Belinfante 改进要解决的问题。

\section{Belinfante--Rosenfeld 改进}

正则 能动量张量 $T^{\mu\nu}_{\text{can}}$ 有两个问题。第一，它一般不对称（见\eqnrefs{eq:ms-poincare-T-asym}）。第二，它不是直接从度规变分得到的，因此与广义相对论中的 Hilbert 能动量张量不一致。Belinfante--Rosenfeld 改进通过添加一个 superpotential 的散度来同时解决这两个问题。

定义 Belinfante 修正
\begin{equation}\label{eq:ms-poincare-B-def}
B^{\lambda\mu\nu}=\frac{1}{2}\big(S^{\lambda\mu\nu}+S^{\mu\nu\lambda}+S^{\nu\mu\lambda}\big),
\end{equation}
改进的能动量张量
\begin{equation}\label{eq:ms-poincare-TB-def}
\boxed{
T_B^{\mu\nu}=T_{\text{can}}^{\mu\nu}+\partial_\lambda B^{\lambda\mu\nu}.
}
\end{equation}

\begin{derivation}[Belinfante 张量的对称性与守恒性]
利用反对称关系\eqnrefs{eq:ms-poincare-T-asym} 与 Belinfante 修正\eqnrefs{eq:ms-poincare-B-def} 的指标性质，分三步验证：(1) $T_B^{\mu\nu}=T_B^{\nu\mu}$（对称性），(2) $\partial_\mu T_B^{\mu\nu}=0$（守恒性），(3) 角动量流可用 $T_B$ 单独表示（自旋吸收）。

先计算 $T_B^{\mu\nu}-T_B^{\nu\mu}$。由定义\eqnrefs{eq:ms-poincare-TB-def}，
\begin{equation*}
T_B^{\mu\nu}-T_B^{\nu\mu}=(T_{\text{can}}^{\mu\nu}-T_{\text{can}}^{\nu\mu})+\partial_\lambda(B^{\lambda\mu\nu}-B^{\lambda\nu\mu}).
\end{equation*}
第一项由\eqnrefs{eq:ms-poincare-T-asym} 给出：$T_{\text{can}}^{\mu\nu}-T_{\text{can}}^{\nu\mu}=-\partial_\lambda S^{\lambda\mu\nu}$。第二项需计算 $B^{\lambda\mu\nu}-B^{\lambda\nu\mu}$。由\eqnrefs{eq:ms-poincare-B-def}：
\begin{equation*}
B^{\lambda\mu\nu}=\frac{1}{2}\big(S^{\lambda\mu\nu}+S^{\mu\nu\lambda}+S^{\nu\mu\lambda}\big),\qquad
B^{\lambda\nu\mu}=\frac{1}{2}\big(S^{\lambda\nu\mu}+S^{\nu\mu\lambda}+S^{\mu\nu\lambda}\big).
\end{equation*}
相减时 $S^{\mu\nu\lambda}$ 和 $S^{\nu\mu\lambda}$ 各自抵消：
\begin{equation*}
B^{\lambda\mu\nu}-B^{\lambda\nu\mu}=\frac{1}{2}\big(S^{\lambda\mu\nu}-S^{\lambda\nu\mu}\big).
\end{equation*}
利用自旋流后两指标反对称性 $S^{\lambda\nu\mu}=-S^{\lambda\mu\nu}$：
\begin{equation*}
B^{\lambda\mu\nu}-B^{\lambda\nu\mu}=\frac{1}{2}\big(S^{\lambda\mu\nu}+S^{\lambda\mu\nu}\big)=S^{\lambda\mu\nu}.
\end{equation*}
因此
\begin{equation*}
T_B^{\mu\nu}-T_B^{\nu\mu}=-\partial_\lambda S^{\lambda\mu\nu}+\partial_\lambda S^{\lambda\mu\nu}=0.
\end{equation*}
Belinfante 张量对称。

守恒性。由\eqnrefs{eq:ms-poincare-TB-def}，$\partial_\mu T_B^{\mu\nu}=\partial_\mu T_{\text{can}}^{\mu\nu}+\partial_\mu\partial_\lambda B^{\lambda\mu\nu}$。第一项为零（正则 守恒\eqnrefs{eq:ms-poincare-T-conservation}）。第二项中 $B^{\lambda\mu\nu}$ 对前两指标反对称：$B^{\lambda\mu\nu}=-B^{\mu\lambda\nu}$。验证：由定义
\begin{equation*}
B^{\mu\lambda\nu}=\frac{1}{2}\big(S^{\mu\lambda\nu}+S^{\lambda\nu\mu}+S^{\nu\lambda\mu}\big),
\end{equation*}
利用 $S^{\mu\lambda\nu}=-S^{\lambda\mu\nu}$、$S^{\lambda\nu\mu}=-S^{\lambda\mu\nu}$（后两指标反对称）、$S^{\nu\lambda\mu}=-S^{\nu\mu\lambda}$，可验证 $B^{\mu\lambda\nu}=-B^{\lambda\mu\nu}$。因此 $\partial_\mu\partial_\lambda$（对称导数）乘 $B^{\lambda\mu\nu}$（前两指标反对称）恒为零：
\begin{equation*}
\partial_\mu\partial_\lambda B^{\lambda\mu\nu}=\frac{1}{2}(\partial_\mu\partial_\lambda-\partial_\lambda\partial_\mu)B^{\lambda\mu\nu}=0.
\end{equation*}
故 $\partial_\mu T_B^{\mu\nu}=0+0=0$。

角动量流改写。用 $T_B$ 改写总"纯轨道"角动量流：
\begin{equation*}
J_B^{\lambda\mu\nu}=x^\mu T_B^{\lambda\nu}-x^\nu T_B^{\lambda\mu}.
\end{equation*}
这个形式的角动量流守恒 $\partial_\lambda J_B^{\lambda\mu\nu}=0$ 等价于原来的 $\partial_\lambda(L^{\lambda\mu\nu}+S^{\lambda\mu\nu})=0$，但自旋信息已被吸收到对称的 $T_B$ 中——不再需要单独的自旋流。这就是 Belinfante 改进的物理意义：把"轨道+自旋"的分解重新打包成一个对称张量的"纯轨道"形式。
\end{derivation}

Belinfante 改进后的 $T_B^{\mu\nu}$ 对称且守恒，与广义相对论中通过度规变分得到的 Hilbert 能动量张量一致。这为下一章把场论推广到曲率时空做好了准备：在曲率时空中，能动量张量必须对称才能与引力耦合，而 Belinfante 张量恰好满足这一要求。

\section*{本章小结}

本章从 Poincaré 群的半直积结构出发，建立了一套完整的"群→代数→表示→场→守恒律"链条。关键步骤和物理含义如下：

\begin{enumerate}
\item \textbf{群结构}：Poincaré 群是 $SO(1,d-1)\ltimes\mathbb{R}^{1,d-1}$ 的半直积。群乘法中 Lorentz 部分作用于平移部分，这是 $[M,P]\neq 0$ 的根源。
\item \textbf{Lie 代数}：统一的半直积 Lie 括号\eqnrefs{eq:ms-poincare-unified-bracket} 一次性给出全部三组交换关系。$[P,P]=0$（平移 Abel），$[M,P]\sim P$（Lorentz 旋转平移），$[M,M]\sim M$（Lorentz 子代数）。
\item \textbf{场表示}：总 Lorentz 生成元 $M=L+\Sigma$ 分为轨道部分 $L$（坐标移动，所有场共有）和内禀部分 $\Sigma$（场分量旋转，由场类型决定）。标量 $\Sigma=0$，矢量 $\Sigma$ 是 $d\times d$ 矩阵，旋量 $\Sigma=\frac{i}{4}[\gamma^\mu,\gamma^\nu]$。
\item \textbf{Noether 定理}：连续准对称性 $\delta\mathcal{L}=\partial_\mu K^\mu$ 给出守恒流 $J^\mu=\Pi_A^\mu\delta\Phi^A-K^\mu$，在壳 满足 $\partial_\mu J^\mu=0$。
\item \textbf{守恒律}：平移对称性给出能动量守恒 $\partial_\mu T^{\mu\nu}=0$，Lorentz 对称性给出角动量守恒 $\partial_\lambda(L^{\lambda\mu\nu}+S^{\lambda\mu\nu})=0$。
\item \textbf{Belinfante 改进}：通过添加 superpotential 散度把 正则 能动量张量对称化，自旋流被吸收进对称的 $T_B^{\mu\nu}$ 中。
\end{enumerate}

下一章把这套结构推广到曲率时空：平直的 $\partial_\mu$ 被协变导数 $D_\mu=\partial_\mu-\frac{i}{2}\omega_{\mu ab}\Sigma^{ab}$ 取代，全局 Poincaré 对称被 微分同胚 协方差 与局域 Lorentz 对称取代，并系统讨论自旋与螺旋度的 Wigner 分类。
